%%%%%%%%%%%%%%%%%%%%%%%%%%%%%%%%%%%%%%%%%%%%%%%%%%%%%%%%%%%%%%
%%%%%%%%%%%%%%%%%% MA-0821 Teoría de módulos %%%%%%%%%%%%%%%%%%%%%%%%%%%%%%
%%%%%%%%%%%%%%%%%%%%%%%%%%%%%%%%%%%%%%%%%%%%%%%%%%%%%%%%%%%%%%
\newpage
\subsection*{MA-0821 Teoría de módulos}
\addcontentsline{toc}{subsection}{MA-0821 Teoría de módulos}

\hypertarget{MA0821}{}
\begin{refsection}
\begin{table}[H]
\centering{
\begin{tabular}{|ll|}
\hline
\textbf{Año:} IV  & \textbf{Requisitos:} Por determinar \\ 
\textbf{Ciclo:} Optativa  & \textbf{Correquisitos:} Ninguno \\ 
\textbf{Tipo de curso:} Teórico & \textbf{Horas presenciales por semana:} Por determinar \\ 
\textbf{Créditos:} 5 & \textbf{Horas de trabajo independiente por semana:} Por determinar \\
\textbf{Virtualidad:} P  &  \\ \hline
\end{tabular}
}
\end{table}

\subsubsection*{Descripción del curso}
Este curso introduce la teoría de módulos como generalización natural de los espacios vectoriales y los grupos abelianos. Se estudian los módulos sobre anillos arbitrarios, con énfasis en los módulos sobre dominios de ideales principales y sus aplicaciones a la clasificación de grupos abelianos finitamente generados y a la forma normal de Jordan. %El curso también cubre nociones de álgebra homológica elemental, incluyendo sucesiones exactas, funtores $\mathrm{Hom}$ y $\otimes$, y propiedades de proyectividad, inyectividad y planicidad.



%\textit{Programa pendiente de elaboración.}

\subsubsection*{Objetivos}

\textbf{General:} Comprender los fundamentos de la teoría de módulos y sus principales aplicaciones algebraicas.

\textbf{Específicos:}

\begin{enumerate}
        \item Comprender la definición y propiedades fundamentales de los módulos sobre anillos, reconociéndolos como marco unificador de estructuras algebraicas previas.
    \item Estudiar submódulos, cocientes, homomorfismos y los teoremas de isomorfismo en el contexto de módulos.
    \item Analizar módulos libres, proyectivos, inyectivos y planos, así como sus caracterizaciones y ejemplos.
    \item Clasificar los módulos finitamente generados sobre dominios de ideales principales y aplicar este resultado a grupos abelianos y formas canónicas de operadores lineales.
    \item Introducir las nociones básicas de álgebra homológica: sucesiones exactas, resoluciones proyectivas e inyectivas, y los funtores derivados $\mathrm{Tor}$ y $\mathrm{Ext}$.
    \item Desarrollar la capacidad de leer y producir demostraciones rigurosas en álgebra abstracta avanzada.
\end{enumerate}


\subsubsection*{Contenidos}

\begin{enumerate}
    \item \textbf{Módulos: definiciones y ejemplos fundamentales (2 semanas)}
    \begin{itemize}
        \item Definición de módulo izquierdo, derecho y bimódulo
        \item Ejemplos: espacios vectoriales, ideales, grupos abelianos, módulos sobre álgebras de grupo
        \item Submódulos y submódulos generados por un conjunto
        \item Módulos cociente y homomorfismos de módulos
        \item Teoremas de isomorfismo para módulos
    \end{itemize}

    \item \textbf{Construcciones básicas (2 semanas)}
    \begin{itemize}
        \item Suma directa interna y externa
        \item Producto directo
        \item Módulos libres y bases
        \item Presentaciones de módulos
        \item Módulo de homomorfismos $\mathrm{Hom}_R(M,N)$
    \end{itemize}

    \item \textbf{Módulos noetherianos y artinianos (2 semanas)}
    \begin{itemize}
        \item Condición de cadena ascendente y descendente
        \item Módulos y anillos noetherianos
        \item Módulos finitamente generados sobre anillos noetherianos
        \item Longitud de un módulo y series de composición
        \item Teorema de Jordan--Hölder
    \end{itemize}

    \item \textbf{Módulos sobre dominios de ideales principales (4 semanas)}
    \begin{itemize}
        \item Submódulos de módulos libres sobre un DIP
        \item Teorema de clasificación de módulos finitamente generados sobre un DIP
        \item Forma invariante (divisores elementales) y forma primaria
        \item Aplicaciones: clasificación de grupos abelianos finitamente generados
        \item Aplicaciones: forma normal de Jordan y forma racional de operadores lineales
    \end{itemize}

    \item \textbf{Módulos proyectivos, inyectivos y planos (2 semanas)}
    \begin{itemize}
        \item Sucesiones exactas cortas y escisión
        \item Módulos proyectivos: definición, caracterizaciones y ejemplos
        \item Módulos inyectivos: definición, envolturas inyectivas y Teorema de Baer
        \item Módulos planos: definición y relación con proyectivos
        \item Lema de Schanuel
    \end{itemize}

    \item \textbf{Producto tensorial (1 semanas)}
    \begin{itemize}
        \item Construcción del producto tensorial $M \otimes_R N$
        \item Propiedades universales y ejemplos
        \item Exactitud del funtor tensorial y módulos planos
        \item Extensión y restricción de escalares
    \end{itemize}

    \item \textbf{Introducción al álgebra homológica (Optativo)(2 semanas)}
    \begin{itemize}
        \item Complejos de cadenas y homología
        \item Resoluciones proyectivas e inyectivas
        \item Funtores derivados: $\mathrm{Ext}^n_R(M,N)$ y $\mathrm{Tor}_n^R(M,N)$
        \item Dimensión proyectiva e inyectiva
        \item Dimensión global de un anillo
    \end{itemize}
\end{enumerate}


\subsubsection*{Metodología}


\noindent \textbf{Lineamientos metodológicos:} La persona docente a cargo de este curso debe respetar las siguientes pautas:
\begin{enumerate}
    \item El curso debe combinar el desarrollo riguroso de la teoría con ejemplos concretos de módulos sobre anillos específicos.
    \item Siempre que resulte natural, se debe enfatizar la conexión de los módulos con estructuras algebraicas ya conocidas por el estudiantado, como espacios vectoriales, grupos abelianos e ideales.
    \item La teoría de módulos sobre dominios de ideales principales debe desarrollarse con detalle, destacando sus aplicaciones a la clasificación de grupos abelianos finitamente generados y a las formas canónicas de operadores lineales.
    \item Las nociones de módulo proyectivo, inyectivo y plano deben presentarse con claridad, ilustrando mediante ejemplos las diferencias entre estas clases.
    \item La introducción al álgebra homológica debe limitarse a los conceptos fundamentales necesarios para un primer curso; no se requiere desarrollar teoría espectral ni métodos homológicos avanzados.
\end{enumerate}

\textbf{Sugerencias metodológicas:} Adicionalmente, se tienen las siguientes recomendaciones para la persona docente a cargo de este curso:
\begin{enumerate}
    \item Utilizar $\mathbb{Z}$, anillos de polinomios y álgebras de grupo como fuentes recurrentes de ejemplos a lo largo del curso.
    \item Presentar el producto tensorial y el funtor $\mathrm{Hom}$ con ejemplos concretos antes de desarrollar sus propiedades generales.
    \item Motivar las nociones de proyectividad e inyectividad mediante su papel en la escisión de sucesiones exactas cortas.
    \item Relacionar el teorema de clasificación sobre dominios de ideales principales con resultados conocidos del álgebra lineal, como la forma normal de Jordan y la forma racional.
    \item Complementar las demostraciones abstractas con ejercicios de cálculo explícito en anillos concretos.
    \item En la evaluación, combinar problemas de demostración con ejercicios de cálculo algebraico directo.
    \item Promover la lectura de secciones seleccionadas de distintas referencias para comparar enfoques y desarrollar autonomía matemática.
    \item Promover la comunicación clara y rigurosa de argumentos e ideas matemáticas de manera oral y escrita.
\end{enumerate}


\nocite{dummit2004abstract}
\nocite{lang2002algebra}
\nocite{rotman2010advanced}
\nocite{weibel1994homological}
\nocite{rotman2009homological}
\nocite{atiyah1969commutative}
\nocite{jacobson1989basic}
\printcursosbibliography

\end{refsection}
